% ============================================================================
% WORKBOOK (MODIFIED) — Section 2.6: Simple Interest
% CCSS 7.RP.A.3 (FL/TX: financial literacy emphasis)
% Practice-focused reinforcement with savings vs. borrowing comparison
% ============================================================================
\section{Simple Interest}
\topicTitle{Simple Interest}

\begin{quickReview}{Simple Interest --- Saving vs.\ Borrowing}
Simple interest uses the formula:
$$I = P \times r \times t$$
\begin{itemize}[leftmargin=*]
    \item $P$ = principal (starting amount), \; $r$ = annual rate (as a decimal), \; $t$ = time in years
    \item \textbf{Saving:} you \emph{earn} interest. Total $= P + I$.
    \item \textbf{Borrowing:} you \emph{pay} interest. Total owed $= P + I$.
\end{itemize}

\medskip
\textbf{Example:} You deposit $\$500$ at $3\%$ for $2$ years.

$I = 500 \times 0.03 \times 2 = \$30$. Total $= 500 + 30 = \$530$. You earned $\$30$.
\end{quickReview}

% ── Warm-Up ──────────────────────────────────────────────────────────────────
\begin{practiceBox}[funGreen]{\faSun~Warm-Up}

\practiceHeader[funGreen]{Find the simple interest and total.}

\prob $P = \$200$, $r = 5\%$, $t = 3$ years. Interest? \answerBlank[3cm]
\answerExplain{$\$30$}{$I = 200 \times 0.05 \times 3 = \$30$.}

\prob $P = \$1{,}000$, $r = 4\%$, $t = 2$ years. Total? \answerBlank[3cm]
\answerExplain{$\$1{,}080$}{$I = 1{,}000 \times 0.04 \times 2 = \$80$. Total $= 1{,}000 + 80 = \$1{,}080$.}

\prob $P = \$750$, $r = 6\%$, $t = 1$ year. Interest? \answerBlank[3cm]
\answerExplain{$\$45$}{$I = 750 \times 0.06 \times 1 = \$45$.}

\prob $P = \$400$, $r = 3.5\%$, $t = 4$ years. Total? \answerBlank[3cm]
\answerExplain{$\$456$}{$I = 400 \times 0.035 \times 4 = \$56$. Total $= 400 + 56 = \$456$.}

\prob $P = \$1{,}500$, $r = 2\%$, $t = 5$ years. Interest? \answerBlank[3cm]
\answerExplain{$\$150$}{$I = 1{,}500 \times 0.02 \times 5 = \$150$.}
\end{practiceBox}

% ── Practice A: Saving Money ─────────────────────────────────────────────────
\begin{practiceBox}{Saving Money}

\practiceHeader{Calculate the interest earned and the total in each savings account.}

\prob You deposit $\$600$ in a savings account at $3\%$ annual interest for $4$ years. How much interest do you earn? What is the total?
\answerBlank[4cm]
\answerExplain{Interest $= \$72$; Total $= \$672$}{$I = 600 \times 0.03 \times 4 = \$72$. Total $= 600 + 72 = \$672$.}

\prob A grandparent puts $\$2{,}000$ in an account at $4.5\%$ for $3$ years for a college fund. How much will be in the account?
\answerBlank[4cm]
\answerExplain{$\$2{,}270$}{$I = 2{,}000 \times 0.045 \times 3 = \$270$. Total $= 2{,}000 + 270 = \$2{,}270$.}

\prob You save $\$350$ at $2.5\%$ for $6$ years. How much interest do you earn?
\answerBlank[3cm]
\answerExplain{$\$52.50$}{$I = 350 \times 0.025 \times 6 = \$52.50$.}

\prob Which earns more interest: $\$800$ at $3\%$ for $5$ years, or $\$1{,}000$ at $2\%$ for $5$ years?
\answerBlank[4cm]
\answerExplain{They are equal: both earn $\$100$}{Option A: $I = 800 \times 0.03 \times 5 = \$120$. Option B: $I = 1{,}000 \times 0.02 \times 5 = \$100$. Option A earns more.}

\end{practiceBox}

% ── Practice B: Borrowing Money ──────────────────────────────────────────────
\begin{practiceBox}[funTeal]{Borrowing Money}

\practiceHeader[funTeal]{Calculate the interest owed and the total repayment.}

\prob You borrow $\$500$ at $8\%$ for $2$ years. How much interest do you pay? What is the total to repay?
\answerBlank[4cm]
\answerExplain{Interest $= \$80$; Repay $= \$580$}{$I = 500 \times 0.08 \times 2 = \$80$. Total $= 500 + 80 = \$580$.}

\prob A family borrows $\$3{,}000$ for a used car at $6\%$ for $3$ years. What is the total repayment?
\answerBlank[4cm]
\answerExplain{$\$3{,}540$}{$I = 3{,}000 \times 0.06 \times 3 = \$540$. Total $= 3{,}000 + 540 = \$3{,}540$.}

\prob You borrow $\$1{,}200$ at $10\%$ for $18$ months. What interest do you owe?
\answerBlank[3cm]
\answerExplain{$\$180$}{$t = 18 \div 12 = 1.5$ years. $I = 1{,}200 \times 0.10 \times 1.5 = \$180$.}

\prob A student borrows $\$2{,}500$ at $5.5\%$ for $4$ years. How much more than the principal will they repay?
\answerBlank[3cm]
\answerExplain{$\$550$ more}{$I = 2{,}500 \times 0.055 \times 4 = \$550$. The extra amount repaid is the interest.}

\end{practiceBox}

% ── Practice C: Saving vs. Borrowing Comparison ─────────────────────────────
\begin{practiceBox}[funOrange]{Saving vs.\ Borrowing}

\practiceHeader[funOrange]{Compare the cost of borrowing to the benefit of saving.}

\prob You can save $\$1{,}000$ at $3\%$ for $3$ years, or borrow $\$1{,}000$ at $9\%$ for $3$ years. Find the interest for each. How much better off are you by saving?
\answerExplain{Save: earn $\$90$; Borrow: pay $\$270$; $\$360$ better off saving}{Saving: $I = 1{,}000 \times 0.03 \times 3 = \$90$ earned. Borrowing: $I = 1{,}000 \times 0.09 \times 3 = \$270$ paid. Difference: $90 + 270 = \$360$ better off.}

\prob Account A offers $2\%$ on savings. Account B offers $4\%$ but charges a $\$25$ annual fee. You deposit $\$500$ for $3$ years. Which gives you more money?
\answerExplain{Account A: $\$530$; Account B: $\$485$; A is better}{A: $I = 500 \times 0.02 \times 3 = \$30$, total $= \$530$. B: $I = 500 \times 0.04 \times 3 = \$60$, fees $= 3 \times 25 = \$75$, total $= 500 + 60 - 75 = \$485$. Account A wins.}

\prob A store offers two payment plans for a $\$400$ TV: (A) Pay $\$400$ now. (B) Pay $\$100$ now and borrow $\$300$ at $12\%$ for $1$ year. What is the total cost of Plan B?
\answerExplain{$\$436$}{Interest on loan: $I = 300 \times 0.12 \times 1 = \$36$. Total: $100 + 300 + 36 = \$436$. Plan B costs $\$36$ more than paying cash.}

\end{practiceBox}

% ── Practice D: True or False ────────────────────────────────────────────────
\begin{practiceBox}[funPurple]{True or False?}

\trueOrFalse{Simple interest grows by the same dollar amount every year.}
\answerExplain{True}{With simple interest, the interest each year is always $P \times r$, which stays constant.}

\trueOrFalse{Borrowing at $6\%$ for $2$ years costs the same interest as borrowing at $3\%$ for $4$ years (same principal).}
\answerExplain{True}{Both give $I = P \times 0.06 \times 2 = P \times 0.12$ and $I = P \times 0.03 \times 4 = P \times 0.12$. The products are equal.}

\trueOrFalse{If you double the principal, the interest also doubles (same rate and time).}
\answerExplain{True}{$I = P \times r \times t$. Doubling $P$ doubles $I$ because $r$ and $t$ stay the same.}

\trueOrFalse{A savings account earning $2\%$ will double your money in $25$ years with simple interest.}
\answerExplain{False}{In $25$ years: $I = P \times 0.02 \times 25 = 0.50P$. Total $= 1.50P$, not $2P$. It takes $50$ years to double at $2\%$ simple interest.}

\end{practiceBox}

% ── Word Problems ────────────────────────────────────────────────────────────
\begin{practiceBox}[funBlue]{\faGlobe~Word Problems}

\wordProblem{Maria deposits $\$800$ at $3.5\%$ simple interest. After how many years will she have earned $\$140$ in interest?}{years}
\answerExplain{$5$ years}{$I = P \times r \times t$. Solve for $t$: $140 = 800 \times 0.035 \times t$, so $t = 140 \div 28 = 5$ years.}

\wordProblem{Jake borrows $\$1{,}500$ at $7\%$ for a summer project. He pays the loan back after $9$ months. How much interest does he owe?}{dollars}
\answerExplain{$\$78.75$}{$t = 9 \div 12 = 0.75$ years. $I = 1{,}500 \times 0.07 \times 0.75 = \$78.75$.}

\end{practiceBox}

% ── Challenge ────────────────────────────────────────────────────────────────
\begin{challengeBox}
\prob Two banks offer different deals on a $\$2{,}000$ deposit. Bank A: $3.5\%$ simple interest. Bank B: $3\%$ simple interest with a $\$20$ sign-up bonus added to the principal. After $5$ years, which bank gives you more money? \answerBlank[5cm]
\answerExplain{Bank A gives $\$2{,}350$; Bank B gives $\$2{,}323$; Bank A is better}{A: $I = 2{,}000 \times 0.035 \times 5 = \$350$, total $= \$2{,}350$. B: Principal $= 2{,}020$, $I = 2{,}020 \times 0.03 \times 5 = \$303$, total $= 2{,}020 + 303 = \$2{,}323$.}

\prob You want to earn $\$200$ in interest from a $\$1{,}600$ deposit in exactly $2$ years. What annual rate do you need? \answerBlank[3cm]
\answerExplain{$6.25\%$}{$200 = 1{,}600 \times r \times 2$. So $200 = 3{,}200r$, and $r = 200 \div 3{,}200 = 0.0625 = 6.25\%$.}
\end{challengeBox}

\encouragement{Understanding interest is a superpower! You now know how money grows when you save --- and what it costs when you borrow!}
